\documentclass[12pt, a4paper]{article}

% --- Настройка русского языка и шрифтов ---
\usepackage[T2A]{fontenc}
\usepackage[utf8]{inputenc}
\usepackage[russian, english]{babel}
\usepackage{amsmath, amssymb}
\usepackage[style=numeric]{biblatex}
\usepackage{hyphenat}       % расширенная поддержка переносов
\hyphenation{
    Рас-смат-ри-ва-ет-ся
    про-из-во-ди-тель-ность
    зна-чи-тель-ных
}

\usepackage{graphicx}
\usepackage{geometry}
\usepackage{multicol}
\usepackage{hyperref}

\usepackage{float}


\addbibresource{bibliography/references.bib}

\geometry{top=2cm, bottom=2cm, left=2.5cm, right=1.5cm}

\begin{document}

% --- Заголовок и автор ---
\begin{center}
    \textbf{\Large Оптимизация при помощи LLM} \\[0.5cm]
\end{center}
    
\section*{Введение}
Задача SAT (задача поиска выполняющего набора для булевой функции в КНФ) является NP-полной. Соответственно, эффективные солверы SAT позволяют эффективно решать задачи комбинаторной оптимизации. Современные CDCL-солверы (Conflict-Driven Clause Learning), такие как Kissat и CaDiCaL, являются state-of-the-art, однако их производительность сильно зависит от встроенных эвристик. Подбор таких эвристик вручную требует значительных временных затрат и богатого опыта.

\section*{Обзор литературы}
В статье \cite{1} впервые упоминается использование LLM для подбора эвристик для SAT солвера. Эксперименты проводятся на основе 12 датасетов, взятых из SAT Competition 2023 и построенных на основе известных задач. Эти же датасеты взяты за основу и в текущей работе. В качестве эвристик предлагаются 9 основных функций солвера EasySAT, который показывает на предложенных датасетах результаты хуже, чем state-of-the-art солверы, но при этом размер EasySAT равен 5 тысячам токенов (в то же время MiniSAT - 25 тысяч, Kissat - 250 тысяч). EasySAT с улучшением эвристик превзошел state-of-the-art солверы на 4 из 12 датасетах. 

В \cite{2} предложена улучшенная структура солвера, с меньшей максимальной длиной функций и с предварительным исследованием предлагаемых эвристик. Также LLM теперь оптимизирует не только сами эвристики, но и промпт для запроса этих эвристик. Рассматривается 10 датасетов из прошлой статьи, и предложенный солвер превосходит state-of-the-art в 8 из 10 датасетах. 

Предыдущие работы сакцентированы на подтверждении концепта поиска релевантных эвристик с помощью LLM для рассматриваемого датасета, но не проверке обобщающих способностей этих эвристик.

\section*{Постановка задачи}
Целью работы является автоматизация и исследование процесса подбора релевантных эвристик с помощью больших языковых моделей (LLM). Задача заключается в минимизации функции потерь $f$ (в нашем случае PAR2) на наборе задач SAT $P$ в пространстве эвристик $H$:
$$h^{*} \in \arg\min_{h \in H} f(h, P)$$

Важной частью работы является проверка обобщающей способности и подбор параметров (стратегия оптимизации, температура LLM, выбор самой LLM) солвера \cite{1}. 

\section*{Методология}
За основу датасетов, на которых проводились эксперименты, взяты датасеты из \cite{1}. Ради уменьшения времени обучения один из датасетов (Zamkeller) случайным образом засэмплен до размера 35 (было 80), второй (cryptography-ascon) взят в полном объеме. Обучение проходило следующим образом: на 60 итераций (итерация состоит из запроса к LLM и вычислении на датасете с обновленными эвристиками). LLM после каждой итерации сообщалось PAR2 бейзлайна и, в качестве фидбэка, PAR2, достигнутый на последней итераций. Далее происходила стадия вычислений на валидационной выборке.

\section*{Текущие результаты и сложности}
На текущий момент реализована инфраструктура для запусков экспериментов из \cite{1} и проведено обучение на ограниченной выборке данных. Исследовано влияние температуры запросов к LLM и выбора модели на результаты солвера. Реализован планнинг вычислений задач, позволяющий оптимизировать вычисления на датасетах (не влияет на метрики).

Красной пунктирной линией отсекается результат бейзлайна. t - температура, func - число оптимизируемых функций. GHC и EA - стратегии оптимизации, детально объяснены в \cite{1}.

\begin{figure}[H]
    \centering
    \includegraphics[width=0.4\linewidth]{images/par2\_vs\_iter.png}
    \caption{DeepSeek v3.1, t=0.9, dataset: cryptography-ascon, strategy: GHC, func: 4}
    \label{fig:placeholder}
\end{figure}

\begin{figure}[H]
    \centering
    \includegraphics[width=0.4\linewidth]{images/par2\_vs\_iter (1).png}
    \caption{DeepSeek v3.1, t=0.8, dataset: cryptography-ascon, strategy: GHC, func: 4}
    \label{fig:placeholder}
\end{figure}

\begin{figure}[H]
    \centering
    \includegraphics[width=0.4\linewidth]{images/par2\_vs\_iter (2).png}
    \caption{DeepSeek v3.1, t=0.6, dataset: cryptography-ascon, strategy: GHC, func: 4}
    \label{fig:placeholder}
\end{figure}

\begin{figure}[H]
    \centering
    \includegraphics[width=0.4\linewidth]{images/par2\_vs\_iter (4).png}
    \caption{gpt-oss-120b, t=0.6, dataset: cryptography-ascon, strategy: GHC, func: 1}
    \label{fig:placeholder}
\end{figure}

\begin{figure}[H]
    \centering
    \includegraphics[width=0.4\linewidth]{images/par2\_vs\_iter (5).png}
    \caption{gpt-oss-120b, t=0.6, dataset: cryptography-ascon, strategy: GHC, func: 4}
    \label{fig:placeholder}
\end{figure}

\begin{figure}[H]
    \centering
    \includegraphics[width=0.4\linewidth]{images/par2\_vs\_iter (6).png}
    \caption{gpt-oss-120b, t=0.6, dataset: cryptography-ascon, strategy: EA, func: 4}
    \label{fig:placeholder}
\end{figure}


% сюда надо картинки с обучением
Увеличение температуры позволяет уменьшить PAR2 солвера, но и порождает некоторые трудности.
В ходе работы выделены следующие проблемы:
\begin{itemize}
    \item Высокая вычислительная сложность обучения из-за необходимости многократных запусков солвера.
    \item При увеличении температуры модель чаще предлагает некомпилируемый код.
\end{itemize}
Для решения этих проблем на данном этапе оптимизация ограничена 4 функциями, урезанными датасетами. Также инжиниринг фидбэк промптов существенно уменьшил число возникающих галлюцинаций.  


\section*{Дальнейшая работа}
Планируется добавление поддержки оптимизации 9 функций, исследование перфоманса более современных моделей, модификация EA стратегии добавлением информации о некотором наборе результатов предыдущих итераций. После обучения достаточного числа солверов будет проведено исследование и сравнение их обобщающих способностей на валидационной выборке.

\printbibliography[heading=bibnumbered]


\end{document}